\documentclass{article}
\usepackage{a4}
\usepackage{amssymb}
\usepackage{graphicx}

\begin{document}
\title{Regularities in electromagnetic decay widths}
\author{Alejandro RIVERO \thanks{EUPT, 
Univ de Zaragoza, Campus de Teruel, 
 44003 Teruel, Spain} \thanks{email: arivero@unizar.es}}

\maketitle

\begin{abstract}
We revisit Sakurai's remark on the regularities of lepton-pair widths for
mesons, extending the panorama to radiative $\to \gamma \gamma$ and 
$\to X \gamma$ decays. The regularities persist, and somehow surprisingly
they seem to relate with Fermi Constant -or Z^0-.
\end{abstract}

Back in 1978, Sakurai took the oportunity of a festscrift  \cite{Sakurai:1978xb} to 
remark how the study of 
vector meson decays $V \to \gamma \to e^+ e^-$, when extended to $J/\Psi$ and $\Upsilon$,
seemed to have confirmed an empirical rule of Yennie about the universality of  $\Gamma (V \to e^+ e^-)$. On other hand a universality in the efective coupling 
$g_{X\gamma\gamma}$ for pseudoscalars $X=\pi,\eta$ is well known popularly and we have recently
noticed that it seems to extend to vector mesons, particularly to $J/\Psi$, when interpretated
as a virtual process for total decay width. And also radiative transitions as $\Sigma^0 \to
\Lambda^0 \gamma$ happen to be in the adequate range. 

Our purpose here is to review all the radiative decays well measured in the modern data tables
in order to establish how significative these regularities are.

\section{All total decay widths}
It is interesting to get some perspective by plotting all the decay rates of subnuclear 
particles. This can be done straightly from the tables provided by the particle data group \cite{}
in its website. One can see three layers corresponding to decays strong, 
electromagnetic or weak, with the neutron being appart
from the rest of weak
decaying particles due to its quark composition.
The area of strong decay is actively pushed forward by experimentalists, 
and one can see how the actual advances hit on one hand the need of 
having a noticeable width in the resonance plot, 
say $\Gamma<E_0$, and on the other hand enough area. We have plotted 
the lines $\Gamma E_0 = 
(1 GeV)^2$, $\Gamma= k E_0$, ($k=1, 0.5, 0.25$), as reference, but we 
excuse ourselves about
plotting the mass against spin dependence famously reported
for this area. 
When particles can not decay strongly -either because there are no other
hadronic mode availble, or because any available mode
should change quark flavours or because asymptotic freedom activates
OZI rule- they give way to electromagnetic decay or weak decay.  

In the weak decay area spectator models and other approaches (even plain dimensional analysis) point to a $\Gamma \sim m^5$ scaling, and we have drawn one such line from the muon. Note that such line in a spectator model should be used to meet the mass value of the decaying quark, not
of the whole particle. 

The intermediate area, electromagnetic or radiative 
decays (we are using both terms as 
exchangeable ones in this context), is the one we are interested. We 
have marked it in figure \ref{fig1} with a
scaling $\Gamma \propto m^3$, partly becayse of the usual $\pi^0$ decay 
calculation,
partly because of dimensional comparation with weak decays: the weak decays
have gained a dimensionful coupling constant (Fermi constant) from
the breaking of electroweak symmetry, and then an aditional $m^2$.

It is interesting to stop a little to think how should the points in the
plot move if electroweak symmetry were restored, playing with the parameters in
the higgs sector, thus altering the value of Fermi constant until a
phase transition is reached and it disappears. It is even possible to use 
mass formulae of mesons and baryions to determine when a given strong
decay becomes available/unavailable and some particle leaves/enter the
electroweak decay zones.

\begin {figure}
%\centering
\includegraphics[width=28cm]{alldecayscolor.eps}
\caption{Log log plot of http:// ,
 $\Gamma(M)$. Both axis have units in MeV }
\label{fig1}
\end {figure}

Also it is a funny try to look for alternative visualizations of the non
strong sector (hint: in modern notation, the strong sector particles are the
ones that carry a mass value between parenthesis as part of its name). It has
been suggested \cite{mcgregor} to use base 137 for the logs in the decay
width, and also\footnote{Amateur suggestion of Yuri Danoyan} to center the 
masses via an arctan(x/M_0) map, adjusting $M_0$ to the value of the proton.

\section{Non weak, non strong decays}

The list of known particles whose main decay mode is neither strong nor 
weak is very short. In order of increasing mass we meet:

-The neutral pion
-$\eta$, although it has already a mix with pions that makes use of the strong
force.
-$Sigma^0$, the only baryon in the party.
-Lower charmonium, this is $J/\Psi$, $\Psi(S2)$, $\eta_c$ (mostly strong 
mixed), and some 1P wave states: $\xi_{c0},\xi_{c1},\xi_{c2}$.
-Lower bottommonium, this is $\Upsilon(S1)$, $\Upsilon(S2)$, $\Upsilon(S3)$
-$Z^0$ and $W^\pm $. Well, why not? They do not decay strongly, and we can
not say strictly that they decay weakly, no more than we can claim that
Death is dead.

If we calculate the reduced decay widths $\beta\def \Gammma/M^3$ for
the above group, only the $Upsilon$ are noticeablely appart (as it
was already apparent in the figure \ref{fig1})


\begin{table}
particle        
\pi 
\eta 
\Sigma^0 
\J/\Psi 
(\Psi(2S)) 
\Upsilon(1S)
W^/pm
Z^0 
$\beta^\frac12$ (MeV) & 
\caption{[Square root of] Reduced full decay widths $\beta=\Gamma/M^3$ for
 EM (generically, non strong non weak) decaying particles}
\end{table}

As I have noted in the introduction, and we will see now in the next
section, the coincidence between $\pi^0$ and $\eta$ increases if we only
consider $\gamma \gamma$ decay.

\section{All the non negligible radiative widths}


%en el plot de aqui, incluir la full decay width de todas y etiquetar
%cada ``columna''


\section{discussion}  






I want to thank K. Illinsky from te PDG by providing a corrected version of the file mass\_-width-2004.csv

\begin{thebibliography}{20}
%\cite{Sakurai:1978xb}
\bibitem{Sakurai:1978xb}
  J.~J.~Sakurai,
  %``Remarkable Regularity In The Lepton - Pair Widths Of Vector Mesons,''
UCLA/78/TEP/20
%\href{http://www.slac.stanford.edu/spires/find/hep/www?r=ucla\%2F78\%2Ftep\%2F20}{SPIRES entry}
\end{thebibliography}
\end{document}
